\section{Estudiar el cambio mismo: identificar primero la fuerza dominante}

Cuando la IA avanza demasiado rápido, seguir las noticias una por una acumula hechos, pero no necesariamente permite formarse un juicio sobre la dirección. Hyung Won Chung elige otro camino: trata la historia de la arquitectura Transformer como un conjunto de experimentos naturales y se pregunta qué diseños eran razonables con las tareas y la capacidad de cómputo de antes, y cuáles perdieron su ventaja al cambiar la escala, los datos y las formas de interacción. El objetivo de la clase no es demostrar que decoder-only siempre es correcto, sino entrenar un método de investigación transferible.

\subsection{Por qué el análisis histórico no es nostalgia}

La portada coloca ``history'' y ``future'' en la misma oración. Al leer las dos diapositivas siguientes conviene captar primero la metodología: el objeto de estudio no es el nombre de algún modelo, sino la trayectoria del cambio. Si se identifican las variables direccionales que gobiernan la trayectoria, las actualizaciones dispersas pueden condensarse en unas pocas hipótesis verificables. Cada comparación de arquitecturas que sigue debe volver a esta pregunta: ¿registra a un ganador estático o explica por qué la ventaja relativa se desplaza con el regime?

\lecturefigure{hyung-slide-001.jpg}{Tema de la conferencia: moldear el futuro de la IA a partir de la historia del Transformer}{Hyung Won Chung official deck, p.~1}

\lecturefigure{hyung-slide-002.jpg}{No basta con seguir los resultados más recientes: hay que estudiar el cambio mismo}{Hyung Won Chung official deck, p.~2}

\teachervoice{00:02:05--00:02:58, Hyung Won enfatiza: cuando la velocidad del cambio supera la capacidad humana de seguirlo, hay que mirar los sistemas anteriores y trazar el camino de «cómo llegamos hasta aquí». Lo central de la clase no es memorizar la tabla de clasificación de 2024, sino abstraer la dirección a partir del camino.}

\begin{warningbox}{Límites de la evidencia de la edición independiente}
Este video es una edición independiente de 36:30 de la lista de reproducción oficial; su contenido coincide con la segunda mitad de la intervención de Hyung Won en la grabación conjunta de la Lecture 27, pero no incluye la Q\&A conjunta posterior. Estas notas conservan solo la prepared talk que aparece efectivamente en el video independiente y no incorporan la discusión de la Q\&A sobre MLE, RLHF o los límites energéticos.
\end{warningbox}

\subsection{Identify, understand, roll forward}

La tercera diapositiva divide la predicción en tres pasos: primero identificar la dominant driving force, luego entender cómo modifica las ventajas y desventajas relativas entre métodos y, por último, extrapolar la tendencia con cautela. La fuerza dominante no es la única fuerza, sino la variable que explica la mayor parte de la direccionalidad con la precisión que se requiere en ese momento; si el regime cambia, hay que volver a modelar. Antes de leer la figura también conviene distinguir entre «describir el pasado» y «predecir el futuro»: lo primero puede acotarse con datos históricos; lo segundo debe declarar a la vez sus condiciones de falla y sus explicaciones alternativas.

\lecturefigure{hyung-slide-003.jpg}{Método de tres pasos para estudiar el cambio: identificar, entender y extrapolar la fuerza dominante}{Hyung Won Chung official deck, p.~3}

\begin{importantbox}{Definición operativa de dominant driving force}
La fuerza impulsora dominante es el factor que basta para explicar los principales cambios de dirección en una escala de tiempo, un presupuesto y una tolerancia al error dados. Permite al investigador ignorar temporalmente los efectos menores a cambio de poder analizar el problema; pero lo que se ignora debe declararse, y cuando la predicción falla, lo primero que hay que revisar son las retroalimentaciones y restricciones omitidas.
\end{importantbox}

\teachervoice{00:03:03--00:04:07, el ponente resume el proceso de investigación como identify, understand, roll forward, y señala que una predicción no necesita una tasa de acierto alta para ser valiosa. Lo importante es que las hipótesis puedan verificarse, no fingir que se conoce el futuro con precisión.}

\subsection{El experimento de la pluma que cae: cuándo puede ignorarse la resistencia del aire}

Cuando se suelta una pluma desde el reposo, en la realidad actúan a la vez la gravedad, la resistencia del aire, la rotación y la deformación del material; aun así, en la clase solo se modela la gravedad, porque en una distancia corta y con una precisión ordinaria es suficientemente dominante. Si se toma como positiva la dirección hacia abajo, puede escribirse
\[
y(t)=y_0+v_0t+\frac{1}{2}gt^2.
\]
donde $y(t)$ es la posición en el instante $t$, $y_0$ es la posición inicial, $v_0$ es la velocidad inicial y $g$ es la aceleración de la gravedad. La condición para que la fórmula sea válida no es que «la resistencia no exista», sino que «la contribución de la resistencia al error sea aceptable para el problema en cuestión».

\lecturefigure{hyung-slide-004.jpg}{El experimento de la pluma que cae: la gravedad ilustra cómo ignorar las fuerzas secundarias}{Hyung Won Chung official deck, p.~4}

\lecturefigure{hyung-slide-006.jpg}{Cuantas más fuerzas dominantes haya y más complejas sean sus interacciones, más difícil es predecir la trayectoria}{Hyung Won Chung official deck, p.~6}

\begin{knowledgebox}{Lectura de la figura: el eje horizontal no es el número de campos observados}
El eje horizontal representa el \textbf{número de fuerzas dominantes en competencia}. Un sistema puede tener muchas características y, aun así, su dirección puede estar determinada por unos pocos mecanismos; a la inversa, con pocas variables pero retroalimentaciones fuertemente acopladas, también puede ser difícil extrapolar. Lo que la figura respalda es un orden de modelado: primero encontrar el factor de mayor direccionalidad y después decidir cuándo incorporar las interacciones.
\end{knowledgebox}

\teachervoice{00:04:40--00:05:57, el ponente menciona explícitamente la resistencia del aire y las interacciones aerodinámicas complejas, pero decide ignorarlas. La nota de clase no es que «la realidad tiene una sola variable», sino «primero se define la precisión requerida y luego se decide qué variables pueden aproximarse».}

\subsection{Resumen de la sección}

Esta sección establece la disciplina de evidencia de toda la clase: la historia no es una lista de modelos, sino un experimento natural; la fuerza dominante no es la única causa, sino una aproximación con un margen de error declarable; la extrapolación no es una garantía, sino la forma de generar el siguiente conjunto de experimentos. A continuación, este método se aplica a compute, inductive bias y la arquitectura Transformer.
